\section{Exact packing under equal relative deadlines}
\label{sec:exact}

The trace experiment uses $d_i=t_i+\delta$ for a common $\delta$. This special
case permits an exact dynamic program. We keep the assumptions visible because
the result does not extend unchanged to arbitrary deadlines or finite service.

\begin{proposition}[Deadline launch times]
\label{prop:deadlines}
Under zero service time and unlimited simultaneous capacity, some optimal
offline schedule launches every batch at an event deadline.
\end{proposition}

\begin{proof}
Take a nonempty feasible batch launched at $\tau$ and let $d_{\min}$ be the
smallest deadline among its events. Feasibility gives $\tau\leq d_{\min}$.
Every assigned event has release at most $\tau$ and deadline at least
$d_{\min}$, so moving the launch to $d_{\min}$ preserves every assignment.
Same-route batches moved to the same time can be merged; different routes may
co-launch under the unlimited-capacity assumption.
\end{proof}

\begin{lemma}[Contiguous block form]
\label{lem:blocks}
For one route with equal relative deadlines, there is an optimum whose batches
are disjoint contiguous blocks in sorted release order.
\end{lemma}

\begin{proof}
Sort releases so $t_1\leq\cdots\leq t_n$, and order batches by launch time.
Suppose $i<j$, event $i$ is assigned to a later batch, and event $j$ to an
earlier batch. Their feasibility gives
\[
 t_i\leq t_j\leq\tau_{\mathrm{early}}\leq
 \tau_{\mathrm{late}}\leq t_i+\delta\leq t_j+\delta.
\]
Exchanging the two assignments is therefore feasible. Repeated exchanges
remove crossings. An unassigned event between the first and last release of a
batch can then be added to that batch: its release is no later than the batch
launch, and its equal-length deadline is no earlier than the first event's
deadline. Thus each batch may be taken as a contiguous block.
\end{proof}

Let $D[j]$ be the maximum number of assigned events among the first $j$
releases of one route, with $D[0]=0$. A feasible block ending at $j$ starts at
some $i\leq j-K_r+1$ and satisfies $t_j-t_i\leq\delta$. Hence

\begin{equation}
\begin{aligned}
D[j]=\max\biggl\{&D[j-1],\\[-2pt]
 &\max_{\substack{i\leq j-K_r+1\\t_j-t_i\leq\delta}}
 \bigl(D[i-1]+j-i+1\bigr)\biggr\}.
\end{aligned}
\label{eq:dp}
\end{equation}

The inner term is $j+\max_i(D[i-1]-i+1)$ over a sliding interval of valid
starts. A grouped sort, route slices, a moving left pointer, and a monotone
deque give an $O(N\log N)$ evaluator including sorting. The frozen evaluator
used for the reported results instead scans the full group array once per route
and finds each left boundary by binary search. Its bound is
$O(NR+\sum_r n_r\log n_r)$ for $R$ routes and $n_r$ events per route, which is
quadratic in the worst case. This affects solver cost, not the returned optimum.
Back-pointers recover a maximum-cardinality witness.

\begin{algorithm}[t]
  \caption{Equal-relative-deadline packing for one route}
  \label{alg:packing}
  \begin{algorithmic}[1]
    \Require sorted integer releases $t_1,\ldots,t_n$, threshold $K$, deadline $\delta$
    \State $D[0]\gets 0$; initialize empty monotone deque $Q$
    \For{$j\gets1$ to $n$}
      \State $D[j]\gets D[j-1]$
      \State $i\gets j-K+1$
      \If{$i\geq1$}
        \State insert $(i,D[i-1]-i+1)$ into $Q$, removing smaller tails
      \EndIf
      \State $\ell\gets\Call{LowerBound}{t,t_j-\delta}$
      \State remove heads whose start index is smaller than $\ell$
      \If{$Q$ is nonempty}
        \State $D[j]\gets\max(D[j],j+Q.\mathrm{head.value})$
      \EndIf
    \EndFor
    \State \Return $D[n]$
  \end{algorithmic}
\end{algorithm}

\paragraph{Exact clock.}
The implementation rounds releases and relative deadlines to integer
nanoseconds and uses inclusive comparisons. The claim is exact on that clock.
It does not rely on a floating tolerance. The implementation agrees with
subset brute force on tiny instances, including route-specific thresholds and
adversarial boundary cases.

\paragraph{Role in this paper.}
For fixed $K$ and $\delta$, the trace oracle is equivalent to selecting the
maximum number of ordered points that can be partitioned into clusters of at
least $K$ points and diameter at most $\delta$. This is the fixed-radius
decision form of one-dimensional $r$-gathering with outliers: related work
often fixes an outlier budget and minimizes radius, whereas we fix the radius
and maximize retained points \cite{aggarwal2010anonymity,akagi2015rgather,
nakano2019rgather}. Classical scheduling also studies compatible batching
under release times and deadlines \cite{barnoy2009throughput}. We use
Equation~\eqref{eq:dp} only as an exact trace oracle and make no algorithmic-priority
claim.
